\documentclass[12pt]{article}
\usepackage[utf8]{inputenc}
\usepackage[T1]{fontenc}
\usepackage{hyperref}
\title{Gepo Language Resources for K-12 Curriculum Development}
\author{Compiled by Claude Sonnet 5.5 (claude-sonnet-5-5) for Andr\'as Kornai}
\date{}
\begin{document}
\maketitle

\section{Introduction}
Gepo (ISO 639-3 \texttt{ygp}; Gepu, Geipo) is a Loloish (Ngwi/Yi) Sino-Tibetan language of Yunnan, China. Wikipedia says it is spoken in six villages of Fumin County, in eastern Luquan County and in adjacent counties. Speaker counts vary by source, from about 62,000 to 100,000. Joshua Project lists 82,000 for Eastern Gepo and 8,400 for Western Gepo, who live in Heqing County. Speakers belong to the Yi nationality and are often called Bai Yi (White Yi). Wikipedia notes that sources use the codes \texttt{ygp} and \texttt{yyz} inconsistently, and classification is debated (Glottolog: Nasu-Gepu; Gao 2017: Central Ngwi).

\section{Primary Educational Resources}
None found in the search results.

\section{Linguistic and Reference Materials}
\begin{itemize}
\item Wikipedia, Gepo language: \url{https://en.wikipedia.org/wiki/Gepo_language}
\item Jiang Jing (2022), \textit{A Study on the Morphological Syntax of Gepo Yi Language}, a grammar monograph listed on Glottolog: \url{https://glottolog.org/resource/reference/id/663389}
\item Glottolog map page for \texttt{gepo1234}: \url{https://glottolog.org/resource/languoid/id/gepo1234.bigmap.html}
\item Numerals data pages, Max Planck Institute for Evolutionary Anthropology: \url{https://lingweb.eva.mpg.de/channumerals/Gepo.htm}
\item Asia Harvest profiles of Eastern and Western Gepo: \url{https://people-groups.asiaharvest.org/China/chinaPeoples/G/GepoEastern.pdf}
\end{itemize}

\section{Religious and Community Resources}
Joshua Project reports that a New Testament was published in 2015 and that no complete Bible exists. Pages: \url{https://joshuaproject.net/index.php/languages/ygp} and \url{https://joshuaproject.net/people_groups/18461} (Eastern Gepo); \url{https://joshuaproject.net/people_groups/18462} (Western Gepo).

\section{Digital and Audio Resources}
Global Recordings Network has a page for the language: \url{https://globalrecordings.net/en/language/ygp}. I did not check what it holds. Glosbe has a ygp page, but I did not check its quality: \url{https://glosbe.com/ygp/de}.

\section{Curriculum Development Context}
Wikipedia lists the Pollard script as the writing system, and I found no standardized orthography or school programme. A curriculum would probably have to be built alongside Yi-language education in China, in consultation with local communities. Since the language is vital, with tens of thousands of speakers, the main need is literacy and teaching materials. No such materials turned up in these results.

\end{document}